\bookchapter{附錄A：新手機開箱“保姆級”流程}

\booksection{開箱驗機“三部曲”：}

\begin{itemize}
  \item \textbf{看外觀}：檢查機身、螢幕有無劃痕、磕碰。
  \item \textbf{對資訊}：核對盒子上的IMEI/序列號和發票上的是否一致。
  \item \textbf{試功能}：撥個電話試試聽筒話筒，拍張照片檢查攝像頭。
\end{itemize}

\booksection{啟用與初始設定“重點提醒”：}

\begin{itemize}
  \item \textbf{網路}：先連上穩定的Wi-Fi。
  \item \textbf{賬號}：登入你的蘋果ID、華為賬號、小米賬號等。
  \item \textbf{基礎設定}：把"查詢我的手機"開啟，設定一個靠譜的鎖屏密碼。
  \item \textbf{eSIM設定}：如果你的手機支援eSIM（如iPhone、部分安卓旗艦），可以在設定中新增eSIM套餐，無需插入實體SIM卡。出國時eSIM尤其方便，可以直接購買當地運營商的eSIM套餐。
  \item \textbf{資料遷移}：從舊手機遷移資料時，確保兩臺手機電量充足。iPhone用"快速開始"，安卓用官方換機App。
\end{itemize}

\booksection{裝App“安全第一”：}

\begin{itemize}
  \item \textbf{首選官方商店}：蘋果用App Store，安卓用手機自帶的應用商店。
  \item \textbf{許可權別亂給}：安裝時彈出的許可權請求，想清楚"這個App真的需要嗎？"
\end{itemize}
